We call the size $K$ down to which we follow the history the \emph{floor}. The \emph{diagonal budget} is the
sum in the following lemma, which controls the cumulative effect of the diagonals on the drifts.
\begin{lema}[Diagonal budget]\label{lema:presupuesto-meso}
Fix $A\ge4C_{\rm geo}$ and a small $\chi>0$. For any integer floor $K_0\le K\le M$, except on an event of
probability $\le C_\chi e^{-c_\chi K}$, we have, up to and including the first-exit board,
\[
 \sum_{\substack{K<k\le M\\\text{regular parents reached}}}
          \frac{(D_{\max}(k)-1)_+}{k^2}
 \le C\{A(\delta+1/n)+\chi+1/K\}.
\]
Here it suffices that $A(\delta+1/n)+\chi+1/K<1/4$. Moreover, for fixed $0<\alpha<1/4$ and any rate
$\lambda>0$, one can take $\delta>0$ sufficiently small and then $K_0,n$ sufficiently large that
\[
 \nu\big(\exists\,K\le k\le M:\ G_k^-,\ D_{\max}(k)\ge\alpha k\big)
 \le C_\lambda e^{-\lambda K},\qquad
 \nu(G_M^-,D_{\max}(M)\ge\alpha M)\le C_\lambda e^{-\lambda M}.
\]
\end{lema}
\begin{proof}
Take the upper endpoints $m$ of the nonempty intervals $[2^jK,\min(2^{j+1}K,M)]$, including $M$ in the last
one, and set $s_m=\lceil Am^2/n+\chi\sqrt{Km}\rceil$ for each of them. Then $s_m/m\le1/4$ and $v=C_{\rm
geo}m^2/(ns_m)\le1/4$. The bound
$4nv^{s_m}=4C_{\rm geo}m^2v^{s_m-1}/s_m\le C m^2e^{-c s_m}$
removes the factor $n$ that comes from the initial board. Summing over the endpoints gives
$C_\chi e^{-c_\chi K}$. In fact, summing $m^2e^{-c\chi\sqrt{Km}}$ over all integers $m\ge K$ gives
the same bound by splitting the sum into bands $2^jK\le m<2^{j+1}K$. If a history reaches an
interval while still regular, it has already reached the upper endpoint of that interval without a
prior exit. Outside the exceptional event, $D_{\max}(m)<s_m$ at that endpoint. Monotonicity and
$s_m-1\le Am^2/n+\chi\sqrt{Km}$ bound the contribution of the interval by $C(Am/n+\chi\sqrt{K/m})$.
We sum the two geometric progressions, including the last, truncated band.

For the second estimate, take $s=\lceil\alpha k\rceil\le k/4$ in \eqref{geompres}. Now $v\le 2C_{\rm
geo}\delta/\alpha$ for large $n$, and $4nv^s\le4C_{\rm geo} k^2v^{s-1}/s$. We choose
$\alpha\log(\alpha/(2C_{\rm geo}\delta))>\lambda$ with strict inequality. The difference absorbs the
polynomial factor and the sum over $k\ge K$, upon increasing $K_0$ if necessary. The estimate at $M$ uses
only the term $k=M$.
\end{proof}

\subsection{Patterns and fibers with an available tracked line}
A fixed pattern $H$ has $s$ row indices, $t$ column indices, a set $S$ of positions that must be available
squares, and a deterministic weight $w$ of bounded norm on each assignment of indices. Interval indicators
may be incorporated into $w$. Let $Z_H$ be the sum over assignments injective within each family, with their
weights and availability indicators, divided by $(k)_s(k)_t$, where $(k)_s=k(k-1)\cdots(k-s+1)$. We take
$K_0$ greater than twice the order of every pattern used.
\begin{lema}[Pattern increments]\label{lema:incrementos-meso}
For a regular parent of size $k$ and $\Delta Z_H=Z_H(k-1)-Z_H(k)$,
\[
 |\mathbb E_\nu(\Delta Z_H\mid\mathcal F^{(k)})|
 \le C_H\frac{1+(D_{\max}(k)-1)_+}{k^2},\qquad
 |\Delta Z_H|\le C_H/k .
\]
The same bounds hold for $F_i(k)/k$, where $i$ is a tracked row and $F_i(k)=\sum_{j\ {\rm
available}}Y_{ij}w_{ij}$, if the step is conditioned on the survival of $i$. They also hold for a tracked
column conditioned to survive. Here $Y_{ij}$ is the availability indicator of the square, and the weights
are fixed and bounded.
\end{lema}
\begin{proof}
By inclusion--exclusion and Lemma~\ref{meso:paso}, the probability that deleting a row and a column affects
a fixed assignment is $(s+t)/k+O_H(k^{-2})$. The exact ratio of the normalizations is $(1-s/k)(1-t/k)$. Thus
the first-order term in the drift cancels. After accounting for the row and column deletions, losing a
square $e\in S$ that is still available requires choosing another square on one of its two diagonals. Each
diagonal offers at most $(D_{\max}(k)-1)_+$ possibilities, each with probability $\le C_q/k^2$. We sum over
$S$ and over assignments, using the absolute value of the weight. For the jump, deleting one row and one
column affects $O_H(k^{s+t-1})$ assignments. When $S\ne\varnothing$, the two deleted diagonals contain at
most $2k$ empty crossings, and fixing a square at a position of $S$ affects at most $O_H(k^{s+t-2})$
assignments. The relative change in normalization is $O_H(1/k)$. The positions of a pattern need not lie on
distinct diagonals.

Now condition only the next step on the survival of $i$. The probability of deleting a column $j$ is
\[
 \frac{\nu(j\ {\rm chosen}\mid\mathcal F^{(k)})-
       \nu((i,j)\ {\rm chosen}\mid\mathcal F^{(k)})}{1-1/k}
   =1/k+O(k^{-2})
\]
uniformly. For a tracked column the families exchange roles and the denominator is $1-1/k+O(k^{-2})$. In
either case the expected loss from the transverse deletion is $F_i/k+O(1/k)$, since there are at most $k$
bounded summands. After division by $k-1$, this cancels the change in normalization up to an error of
$O(k^{-2})$. For each square of the fiber, the additional probability of losing it through a diagonal is at
most $2C_q(D_{\max}(k)-1)_+/k^2$ divided by the probability that the tracked line survives this step, which
is uniformly bounded away from zero. We sum over the fiber and divide by $k-1$. Each diagonal meets the
fiber in at most one square, so the normalized jump is $O(1/k)$.
\end{proof}
